% ==============================================================================
% SECTION 1: INTRODUCTION
% ==============================================================================
% Extensión: 400–600 palabras
% Estructura embudo: contexto general → estado del arte → brecha de conocimiento → objetivo
% Tiempo: presente para hechos generales, pasado para estudios previos
% Incluye: citas con natbib, ecuaciones inline, números formateados (\num{}),
% párrafo guía con referencias cruzadas (\Cref)

\section{Introducción}\label{sec:intro}

La bioinformática contemporánea depende de workflows compuestos por múltiples
herramientas, formatos de archivo y decisiones analíticas que rara vez siguen
una trayectoria completamente lineal. Aunque los pipelines tradicionales han
permitido estandarizar tareas como alineamiento, cuantificación, anotación o
modelado, su ejecución sigue exigiendo intervención humana para adaptar
parámetros, resolver errores, encadenar software heterogéneo y reinterpretar
resultados intermedios. Esta fricción se vuelve más evidente en contextos como
la transcriptómica, las ómicas de célula única o la farmacometría, donde la
complejidad técnica crece al mismo ritmo que el volumen de datos.

En ese escenario, los modelos de lenguaje de gran escala (LLMs) han comenzado a
desempeñar un papel distinto al de asistentes conversacionales genéricos.
Cuando se integran en sistemas capaces de planificar, invocar herramientas,
ejecutar código y corregir errores a partir de la retroalimentación del
entorno, los LLMs pasan de ser interfaces de \textit{prompting} a componentes
de arquitecturas agénticas orientadas a la automatización
científica~\citep{42042548,41912700}. Esta evolución ha favorecido agentes que
operan sobre notebooks, APIs, bases de datos, software bioinformático y
servicios web.

Sin embargo, el interés creciente por estos sistemas no ha venido acompañado de
un marco sintético claro para describir cómo están siendo diseñados, evaluados
y limitados dentro del dominio bioinformático. La literatura reciente incluye
agentes recursivos, sistemas multi-agente, asistentes con \textit{grounding} y
plataformas con observabilidad, pero persisten preguntas sobre
reproducibilidad, coste computacional, trazabilidad, control humano e
interoperabilidad. El campo avanza rápido en prototipos y demostraciones, pero
todavía carece de una visión ordenada de sus arquitecturas dominantes y de sus
condiciones de validez científica.

Esta mini-review examina 15 estudios publicados en 2025 y 2026 sobre el uso de
LLMs y sistemas basados en agentes para la automatización de workflows en
bioinformática. La Sección~\ref{sec:methods} describe la estrategia de búsqueda
y los criterios de selección; la Sección~\ref{sec:synthesis} organiza la
evidencia en cuatro ejes temáticos; la Sección~\ref{sec:discussion} discute los
hallazgos, limitaciones y vacíos de la literatura; y la
Sección~\ref{sec:conclusions} presenta las conclusiones y perspectivas futuras
del área.
